\scanpage{12}
\deleted{$I(U)\otimes_k I(\Lambda)$ le sont.} D'où
\diagram{\deleted{\begin{tikzcd}[ampersand replacement=\&]
I(U)\otimes_k I(\Lambda)\arrow[r] \& 1+I(U)\otimes_k\\
U\otimes_k I(\Lambda) \&
\end{tikzcd}}}
Alors $U\otimes\deleted{\unsure{\Lambda}} I(\Lambda)$ est de puiss. n:ième
nulle, d'où
\diagram{\begin{tikzcd}[ampersand replacement=\&]
U\otimes I(\Lambda)\arrow[r,shift left,"\exp"] \& 1+U\otimes I(\Lambda)\arrow[l,shift left,"\log"]
\end{tikzcd}}
[et si $I(U)$ est \unsure{l'idéal} d'une
augmentation de $U$, on a
\diagram{\begin{tikzcd}[ampersand replacement=\&]
I(U)\otimes I(\Lambda)\arrow[r,shift left,"\exp"] \& 1+I(U)\otimes I(\Lambda)\arrow[l,shift left,"\log"]
\end{tikzcd}}
\hfill$]$

Supposons maintenant que $U$ soit une
hyperalgèbre ordinaire. Alors la correspondance
précédente définit une corr. biunivoque
entre la partie de $\deleted{I(U)\otimes I(\Lambda)}\deleted{\unsure{I(\Lambda\otimes_k U)|I(U\otimes\Lambda)\otimes I(\Lambda)}}\inserted{(U\otimes I(\Lambda))\cap I(U\otimes\Lambda)\unsure{=|\subset} I(U)\otimes I(\unsure{\Lambda|U})}$ formée
des éléments primitifs \deleted{\unsure{sont}} de $U\otimes_k\Lambda$, et
la partie de $1+I(U)\otimes I(\Lambda)$ $\deleted{\subset 1+I(U\otimes\Lambda)}$
formée des élts multiplicatifs. Ce
sont ainsi resp$^{\mathrm t}$. les primitifs de
$U\otimes_k\Lambda$ qui donnent $0$ par l'augmentation
$U\otimes_k\Lambda\longrightarrow U$, et les élémts multiplicatifs de
$U\otimes_k\Lambda$ qui donnent $1$ par ladite augmentation
% archive sheet number: 11

